\documentclass[12pt]{report}  % ou article si tu veux moins structuré
\usepackage[utf8]{inputenc}
\usepackage[T1]{fontenc}
\usepackage[french]{babel}
\usepackage{amsmath, amssymb}
\usepackage{graphicx}
\usepackage{geometry}
\usepackage{titlesec}
\usepackage{xcolor}
\usepackage{fancyhdr}
\usepackage{ulem} % pour souligner le titre

\geometry{a4paper, margin=2.5cm}

% Pour personnaliser les titres
\titleformat{\chapter}[hang]
  {\normalfont\Huge\bfseries\centering}{}{0pt}{\uline}  % souligné + centré

\titleformat{\section}
  {\normalfont\Large\bfseries}{}{0pt}{\uline}

\title{
  \vspace{2cm}
  \textbf{\uline{\Huge Rapport machine learning}}\\[1cm]
}
\author{\Large Clément Berthoud, Christophe Boshra}
\date{6 juin 2025}


% Début du contenu du doc 
\begin{document}

\maketitle

\tableofcontents
\newpage

% ====================
\chapter{Introduction}
\vspace{4cm}
% ← espace vide pour écrire à la main ou compléter plus tard

% ====================
\chapter{Objectifs}
\section{Partie 1}
\vspace{0cm}

Bonjour à toutes et tous 


% ====================
\chapter{Méthodologie}
\section{Collecte des données}
\vspace{4cm}

\section{Analyse}
\vspace{4cm}

% ====================
\chapter{Résultats}
\vspace{4cm}

% ====================
\chapter{Conclusion}
\vspace{4cm}

\end{document}